\ifdefined\gkver\else\def\gkver{student}\fi
\documentclass[\gkver]{gaokaozhenti}
\begin{document}

\chapter{2008年海南}

\begin{enumerate}
\item
%% number: 1
%% paperName: 2008年海南高考第 1 题 3 分
%% typeId: 1
%% type: 单选题
%% chapter: 电磁感应
%% point: 电磁感应现象
%% method: 无
%% score: 3
%% degree: 650
%% duplicateId: 0
%% body:
法拉第通过静心设计的一系列试验，发现了电磁感应定律，将历史上认为各自独立的学科“电学”与“磁学”联系起来。在下面几个典型的实验设计思想中，所作的推论后来被实验否定的是\xzanswer{A}

\fourchoices[answer=A]
{既然磁铁可使近旁的铁块带磁，静电荷可使近旁的导体表面感应出电荷，那么静止导线上的稳恒电流也可在近旁静止的线圈中感应出电流}
{既然磁铁可在近旁运动的导体中感应出电动势，那么稳恒电流也可在近旁运动的线圈中感应出电流}
{既然运动的磁铁可在近旁静止的线圈中感应出电流，那么静止的磁铁也可在近旁运动的导体中感应出电动势}
{既然运动的磁铁可在近旁的导体中感应出电动势，那么运动导线上的稳恒电流也可在近旁的线圈中感应出电流}

%% answer: A
%% memo:
\memoanswer{
对 A 选项，静止的导线上的稳恒电流附近产生稳定的磁场，通过旁边静止的线圈不会产生感应电流，A 被否定；稳恒电流周围的稳定磁场是非匀强磁场，运动的线圈可能会产生感应电流，B 符合事实；静止的磁铁周围存在稳定的磁场，旁边运动的导体棒会产生感应电动势，C 符合；运动的导线上的稳恒电流周围产生运动的磁场，即周围磁场变化，在旁边的线圈中产生感应电流，D 符合。
}

\item
%% number: 2
%% paperName: 2008年海南高考第 2 题 3 分
%% typeId: 1
%% type: 单选题
%% chapter: 力物体的平衡
%% point: 物体系平衡
%% method: 整体与隔离
%% score: 3
%% degree: 650
%% duplicateId: 0
%% body:
如图，质量为 $M$ 的楔形物块静置在水平地面上，其斜面的倾角为 $\theta$。斜面上有一质量为 $m$ 的小物块，小物块与斜面之间存在摩擦。用恒力 $F$ 沿斜面向上拉小物块，使之匀速上滑。在小物块运动的过程中，楔形物块始终保持静止。地面对楔形物块的支持力为\xzanswer{D}

\onepicture[w=4.6cm]{figs/02.png}

\fourchoices[answer=D]
{$(M+m)g$}
{$(M+m)g-F$}
{$(M+m)g+F\sin\theta$}
{$(M+m)g-F\sin\theta$}

%% answer: D
%% memo:
\memoanswer{
本题可用整体法的牛顿第二定律解题，竖直方向由平衡条件：$F\sin\theta+N=mg+Mg$，则 $N=mg+Mg-F\sin\theta$。
}

\item
%% number: 3
%% paperName: 2008年海南高考第 3 题 3 分
%% typeId: 1
%% type: 单选题
%% chapter: 曲线运动
%% point: 竖直平面圆周运动
%% method: 无
%% score: 3
%% degree: 650
%% duplicateId: 0
%% body:
如图，一轻绳的一端系在固定粗糙斜面上的 O 点，另一端系一小球。给小球一足够大的初速度，使小球在斜面上做圆周运动。在此过程中，\xzanswer{C}

\onepicture[w=3.8cm]{figs/03.png}

\fourchoices[answer=C]
{小球的机械能守恒}
{重力对小球不做功}
{绳的张力对小球不做功}
{在任何一段时间内，小球克服摩擦力所做的功总是等于小球动能的减少}

%% answer: C
%% memo:
\memoanswer{
斜面粗糙，小球受到重力、支持力、摩擦力、绳子拉力，由于除重力做功外，摩擦力做负功，机械能减少，A、B 错；绳子张力总是与运动方向垂直，故不做功，C 对；小球动能的变化等于合外力做功，即重力与摩擦力做功，D 错。
}

\item
%% number: 4
%% paperName: 2008年海南高考第 4 题 3 分
%% typeId: 1
%% type: 单选题
%% chapter: 电场
%% point: 电势能电势差电场力做功
%% method: 无
%% score: 3
%% degree: 650
%% duplicateId: 0
%% body:
静电场中，带电粒子在电场力作用下从电势为 $\varphi_{a}$ 的 a 点运动至电势为 $\varphi_{b}$ 的 b 点。若带电粒子在 a、b 两点的速率分别为 $v_{a}$、$v_{b}$，不计重力，则带电粒子的比荷 $q/m$ 为\xzanswer{C}

\fourchoices[answer=C]
{$\frac{v_{a}^{2}-v_{b}^{2}}{\varphi_{b}-\varphi_{a}}$}
{$\frac{v_{b}^{2}-v_{a}^{2}}{\varphi_{b}-\varphi_{a}}$}
{$\frac{v_{a}^{2}-v_{b}^{2}}{2(\varphi_{b}-\varphi_{a})}$}
{$\frac{v_{b}^{2}-v_{a}^{2}}{2(\varphi_{b}-\varphi_{a})}$}

%% answer: C
%% memo:
\memoanswer{
由电势差公式以及动能定理：$W=qU_{ab}=q(\varphi_{a}-\varphi_{b})=\frac{1}{2}m(v_{b}^{2}-v_{a}^{2})$，可得比荷为 $\frac{q}{m}=\frac{v_{b}^{2}-v_{a}^{2}}{2(\varphi_{a}-\varphi_{b})}$。
}

\item
%% number: 5
%% paperName: 2008年海南高考第 5 题 3 分
%% typeId: 1
%% type: 单选题
%% chapter: 电场
%% point: 电场中的图象
%% method: 类比
%% score: 3
%% degree: 450
%% duplicateId: 0
%% body:
质子和中子是由更基本的粒子即所谓“夸克”组成的。两个强作用电荷相反（类似于正负电荷）的夸克在距离很近时几乎没有相互作用（称为“渐近自由”）；在距离较远时，它们之间就会出现很强的引力（导致所谓“夸克禁闭”）。作为一个简单的模型，设这样的两夸克之间的相互作用力 $F$ 与它们之间的距离 $r$ 的关系为：

$F=\left\{\begin{array}{ll}0, & 0<r<r_{1}\\ -F_{0}, & r_{1}\leqslant r\leqslant r_{2}\\ 0, & r>r_{2}\end{array}\right.$

式中 $F_{0}$ 为大于零的常量，负号表示引力。用 $U$ 表示夸克间的势能，令 $U_{0}=F_{0}(r_{2}-r_{1})$，取无穷远为势能零点。下列 $U-r$ 图示中正确的是\xzanswer{B}

\fourchoices[answer=B, ispicture=true, h=1.8cm]
{figs/05a.png}
{figs/05b.png}
{figs/05c.png}
{figs/05d.png}

%% answer: B
%% memo:
\memoanswer{
从无穷远处电势为零开始到 $r=r_{2}$ 位置，势能恒定为零；在 $r=r_{2}$ 到 $r=r_{1}$ 过程中，恒定引力做正功，势能逐渐均匀减小，即势能为负值且越来越小，此部分图像为 A、B 选项中所示；$r<r_{1}$ 之后势能不变，恒定为 $-U_{0}$，由引力做功等于势能减小量，故 $U_{0}=F_{0}(r_{2}-r_{1})$。
}

\item
%% number: 6
%% paperName: 2008年海南高考第 6 题 3 分
%% typeId: 1
%% type: 单选题
%% chapter: 电场
%% point: 等势面
%% method: 无
%% score: 3
%% degree: 250
%% duplicateId: 0
%% body:
匀强电场中有 a、b、c 三点。在以它们为顶点的三角形中，$\angle a=30^{\circ}$、$\angle c=90^{\circ}$，电场方向与三角形所在平面平行。已知 a、b 和 c 点的电势分别为 $(2-\sqrt{3})\UV$、$(2+\sqrt{3})\UV$ 和 $2\UV$。该三角形的外接圆上最低、最高电势分别为\xzanswer{B}

\onepicture[w=3.8cm]{figs/06.png}

\fourchoices[answer=B]
{$(2-\sqrt{3})\UV$、$(2+\sqrt{3})\UV$}
{$0\UV$、$4\UV$}
{$(2-\frac{4\sqrt{3}}{3})\UV$、$(2+\frac{4\sqrt{3}}{3})\UV$}
{$0\UV$、$\sqrt{3}\UV$}

%% answer: B
%% memo:
\memoanswer{
如图，匀强电场的电场线与等势面平行等间距排列，且电场线与等势面处处垂直，沿着电场线方向电势均匀降落。取 ab 的中点 O，即为三角形的外接圆的圆心，且该点电势为 $2\UV$，故 Oc 为等势面，MN 为电场线，方向为 MN 方向，$U_{OP}=U_{Oa}=\sqrt{3}\UV$，$U_{ON}:U_{OP}=2:\sqrt{3}$，故 $U_{ON}=2\UV$，N 点电势为零，为最小电势点；同理 M 点电势为 $4\UV$，为最大电势点。

\onepicture[w=3.6cm]{figs/06_an.png}
}

\item
%% number: 7
%% paperName: 2008年海南高考第 7 题 4 分
%% typeId: 2
%% type: 多选题
%% chapter: 交流电
%% point: 变压器
%% method: 无
%% score: 4
%% degree: 550
%% duplicateId: 0
%% body:
如图，理想变压器原副线圈匝数之比为 $4:1$。原线圈接入一电压为 $u=U_{0}\sin\omega t$ 的交流电源，副线圈接一个 $R=27.5\UO$ 的负载电阻。若 $U_{0}=220\sqrt{2}\UV$，$\omega=100\pi\UHz$，则下述结论正确的是\xzanswer{AC}

\onepicture[w=4.5cm]{figs/07.png}

\fourchoices[answer=AC]
{副线圈中电压表的读数为 $55\UV$}
{副线圈中输出交流电的周期为 $\frac{1}{100\pi}\Us$}
{原线圈中电流表的读数为 $0.5\UA$}
{原线圈中的输入功率为 $110\sqrt{2}\UW$}

%% answer: AC
%% memo:
\memoanswer{
原线圈电压有效值 $U_{1}=220\UV$，由电压比等于匝数比可得副线圈电压 $U_{2}=55\UV$，A 对；电阻 $R$ 上的电流为 $2\UA$，由原副线圈电流比等于匝数的反比，可得电流表示数为 $0.5\UA$，C 对；输入功率为 $P=220\times0.5\UW=110\UW$，D 错；周期 $T=\frac{2\pi}{\omega}=0.02\Us$，B 错。
}

\item
%% number: 8
%% paperName: 2008年海南高考第 8 题 4 分
%% typeId: 2
%% type: 多选题
%% chapter: 直线运动
%% point: 运动图象
%% method: 无
%% score: 4
%% degree: 450
%% duplicateId: 0
%% body:
$t=0$ 时，甲乙两汽车从相距 $70\Ukm$ 的两地开始相向行驶，它们的 $v-t$ 图象如图所示。忽略汽车掉头所需时间。下列对汽车运动状况的描述正确的是\xzanswer{BC}

\onepicture[w=4.2cm]{figs/08.png}

\fourchoices[answer=BC]
{在第 1 小时末，乙车改变运动方向}
{在第 2 小时末，甲乙两车相距 $10\Ukm$}
{在前 4 小时内，乙车运动加速度的大小总比甲车的大}
{在第 4 小时末，甲乙两车相遇}

%% answer: BC
%% memo:
\memoanswer{
速度图像在 $t$ 轴下方的均为反方向运动，故 2h 末乙车改变运动方向，A 错；2h 末从图像围成的面积可知乙车运动位移为 $30\Ukm$，甲车位移为 $30\Ukm$，相向运动，此时两车相距 $70\Ukm-30\Ukm-30\Ukm=10\Ukm$，B 对；从图像的斜率看，斜率大加速度大，故乙车加速度在 4h 内一直比甲车加速度大，C 对；4h 末，甲车运动位移 $120\Ukm$，乙车运动位移 $30\Ukm$，两车原来相距 $70\Ukm$，故此时两车还相距 $20\Ukm$，D 错。
}

\item
%% number: 9
%% paperName: 2008年海南高考第 9 题 4 分
%% typeId: 2
%% type: 多选题
%% chapter: 运动和力
%% point: 牛顿第二定律
%% method: 无
%% score: 4
%% degree: 450
%% duplicateId: 0
%% body:
如图，水平地面上有一楔形物体 b，b 的斜面上有一小物块 a；a 与 b 之间、b 与地面之间均存在摩擦。已知楔形物体 b 静止时，a 静止在 b 的斜面上。现给 a 和 b 一个共同的向左的初速度，与 a 和 b 都静止时相比，此时可能\xzanswer{BC}

\onepicture[w=4.6cm]{figs/09.png}

\fourchoices[answer=BC]
{a 与 b 之间的压力减少，且 a 相对 b 向下滑动}
{a 与 b 之间的压力增大，且 a 相对 b 向上滑动}
{a 与 b 之间的压力增大，且 a 相对 b 静止不动}
{b 与地面之间的压力不变，且 a 相对 b 向上滑动}

%% answer: BC
%% memo:
\memoanswer{
依题意，若两物体依然相对静止，则 a 的加速度一定水平向右，如图将加速度分解为垂直斜面与平行于斜面，则垂直斜面方向，$N-mg\cos\theta=ma_{y}$，即支持力 $N$ 大于 $mg\cos\theta$，与都静止时比较，a 与 b 间的压力增大；沿着斜面方向，若加速度 $a$ 过大，则摩擦力可能沿着斜面向下，即 a 物块可能相对 b 有向上滑动趋势，甚至相对向上滑动，故 A 错，B、C 正确；对系统整体，在竖直方向，若物块 a 相对 b 向上滑动，则 a 还具有向上的分加速度，由整体的牛顿第二定律可知，系统处于超重状态，b 与地面之间的压力将大于两物体重力之和，D 错。

\onepicture[w=3.2cm]{figs/09_an.png}
}

\item
%% number: 10
%% paperName: 2008年海南高考第 10 题 4 分
%% typeId: 2
%% type: 多选题
%% chapter: 电磁感应
%% point: 切割磁感线电动势
%% method: 无
%% score: 4
%% degree: 450
%% duplicateId: 0
%% body:
一航天飞机下有一细金属杆，杆指向地心。若仅考虑地磁场的影响，则当航天飞机位于赤道上空\xzanswer{AD}

\fourchoices[answer=AD]
{由东向西水平飞行时，金属杆中感应电动势的方向一定由上向下}
{由西向东水平飞行时，金属杆中感应电动势的方向一定由上向下}
{沿经过地磁极的那条经线由南向北水平飞行时，金属杆中感应电动势的方向一定由下向上}
{沿经过地磁极的那条经线由北向南水平飞行时，金属杆中一定没有感应电动势}

%% answer: AD
%% memo:
\memoanswer{
如图，设观察方向为面向北方，左西右东，则地磁场方向平行赤道表面向北。若飞机由东向西飞行时，由右手定则可判断出电动势方向为由上向下；若飞机由西向东飞行时，由右手定则可判断出电动势方向为由下向上，A 对 B 错；沿着经过地磁极的那条经线运动时，速度方向平行于磁场，金属杆中一定没有感应电动势，C 错 D 对。

\onepicture[w=5.8cm]{figs/10_an.png}
}

\item
%% number: 11
%% paperName: 2008年海南高考第 11 题 4 分
%% typeId: 3
%% type: 填空题
%% chapter: 电路
%% point: 传感器
%% method: 无
%% score: 4
%% degree: 950
%% duplicateId: 0
%% body:
当光照射到光敏电阻上时，光敏电阻的阻值\tkanswer[1.5cm]{变小}（填“变大”、“不变”或“变小”）。半导体热敏电阻是利用半导体材料的电阻率随\tkanswer[1.5cm]{温度}变化而改变的特性制成的。

%% answer:
\jdanswer{
变小；温度
}

%% memo:
\memoanswer{
光敏电阻和热敏电阻均为半导体材料的电阻，半导体材料的电阻率随温度升高而减小。
}

\item
%% number: 12
%% paperName: 2008年海南高考第 12 题 4 分
%% typeId: 3
%% type: 填空题
%% chapter: 曲线运动
%% point: 万有引力定律
%% method: 无
%% score: 4
%% degree: 650
%% duplicateId: 0
%% body:
一探月卫星在地月转移轨道上运行，某一时刻正好处于地心和月心的连线上，卫星在此处所受地球引力与月球引力之比为 $4:1$。已知地球与月球的质量之比约为 $81:1$，则该处到地心与到月心的距离之比约为\tkanswer[2cm]{9:2}。

%% answer:
\jdanswer{
$9:2$
}

%% memo:
\memoanswer{
由万有引力定律，卫星受到地球和月球的万有引力分别为 $F_{\text{地}}=G\frac{M_{\text{地}}m}{R_{\text{地}}^{2}}$，$F_{\text{月}}=G\frac{M_{\text{月}}m}{R_{\text{月}}^{2}}$，代入题目给定的数据可得 $R_{\text{地}}:R_{\text{月}}=9:2$。
}

\item
%% number: 13
%% paperName: 2008年海南高考第 13 题 4 分
%% typeId: 4
%% type: 实验题
%% chapter: 直线运动
%% point: 游标卡尺和螺旋测微仪
%% method: 无
%% score: 4
%% degree: 850
%% duplicateId: 0
%% body:
某同学用螺旋测微器测量一金属丝的直径，测微器的示数如图所示，该金属丝直径的测量值为\tkanswer[2cm]{2.793}$\Umm$。

\onepicture[w=3.6cm, align=right]{figs/13.png}

%% answer:
\jdanswer{
2.793（2.791～2.795 均可）
}

%% memo:
\memoanswer{
测量值读数为 $2.5\Umm+0.01\Umm\times29.3=2.793\Umm$。
}

\item
%% number: 14
%% paperName: 2008年海南高考第 14 题 11 分
%% typeId: 4
%% type: 实验题
%% chapter: 电路
%% point: 电表改装
%% method: 无
%% score: 11
%% degree: 450
%% duplicateId: 0
%% body:
一毫安表头 mA 满偏电流为 $9.90\UmA$，内阻约为 $300\UO$。要求将此毫安表头改装成量程为 $1\UA$ 的电流表，其电路原理如图所示。图中，A 是量程为 $2\UA$ 的标准电流表，$R_{0}$ 为电阻箱，$R$ 为滑动变阻器，S 为开关，E 为电源。

\begin{enumerate}
	\item
	完善下列实验步骤：
	
	①将虚线框内的实物图按电路原理图连线；
	
	②将滑动变阻器的滑动头调至\tkanswer[1cm]{b}端（填“a”或“b”），电阻箱 $R_{0}$ 的阻值调至零；
	
	③合上开关；
	
	④调节滑动变阻器的滑动头，增大回路中的电流，使标准电流表读数为 $1\UA$；
	
	⑤调节电阻箱 $R_{0}$ 的阻值，使毫安表指针接近满偏，此时标准电流表的读数会\tkanswer[1.2cm]{减小}（填“增大”、“减小”或“不变”）；
	
	⑥多次重复步骤④⑤，直至标准电流表的读数为\tkanswer[1.2cm]{$1\UA$}，同时毫安表指针满偏。
	
	\item
	回答下列问题：
	
	①在完成全部实验步骤后，电阻箱使用阻值的读数为 $3.1\UO$，由此可知毫安表头的内阻为\tkanswer[1.5cm]{$310\UO$}。
	
	②用改装成的电流表测量某一电路中的电流，电流表指针半偏，此时流过电阻箱的电流为\tkanswer[1.8cm]{0.495}$\UA$。
	
	③对于按照以上步骤改装后的电流表，写出一个可能影响它的准确程度的因素：

	\tkanswer[8cm]{}。
\end{enumerate}

\twopicture[align=right, showlabel=false, wA=3.8cm, wB=5.2cm]
{figs/14a.png}
{figs/14b.png}

%% answer:
\jdanswer{
\begin{enumerate}
	\item
	① 连线如图；② b；⑤ 减小；⑥ $1\UA$。

	\onepicture[w=3.2cm]{figs/14_an.jpg}

	\item
	① $310\UO$；② $0.495\UA$（$0.494\sim0.496\UA$ 均可）；③ 例如：电阻箱和滑动变阻器的阻值不能连续变化；标准表和毫安表头的读数误差；电表指针偏转和实际电流的大小不成正比等等。
\end{enumerate}
}

%% memo:
\memoanswer{
（1）将滑动变阻器的滑动头调至 b 点是为了保护毫安表；在调节电阻箱时，由于接入回路的电阻增加，故电流表电流会减小；为了能成功改装为 $1\UA$ 量程的电流表，需要通过调节滑动变阻器使标准电流表示数回到 $1\UA$。

（2）电阻箱的电阻读数为 $3.1\UO$，此时毫安表满偏，电阻箱上的电流为 $1000\UmA-9.9\UmA=990.1\UmA$，根据并联电阻值与电流值成反比，可计算毫安表头的内阻为 $310\UO$；若毫安表半偏，电阻箱的电流也减半为 $495\UmA=0.495\UA$；本实验可能因电阻箱和滑动变阻器的阻值不能连续变化，导致标准电流表的示数无法精准显示 $1\UA$，还有电表本身读数上的误差，以及电表指针偏转本身不是随电流均匀偏转等，都可能影响改装后电流表的准确程度。
}

\item
%% number: 15
%% paperName: 2008年海南高考第 15 题 8 分
%% typeId: 6
%% type: 计算题
%% chapter: 运动和力
%% point: 加速减速模型
%% method: 无
%% score: 8
%% degree: 550
%% duplicateId: 0
%% body:
科研人员乘气球进行科学考察。气球、座舱、压舱物和科研人员的总质量为 $990\Ukg$。气球在空中停留一段时间后，发现气球漏气而下降，及时堵住。堵住时气球下降速度为 $1\Ums$，且做匀加速运动，$4\Us$ 内下降了 $12\Um$。为使气球安全着陆，向舱外缓慢抛出一定的压舱物。此后发现气球做匀减速运动，下降速度在 $5\Umin$ 内减少 $3\Ums$。若空气阻力和泄漏气体的质量均可忽略，重力加速度 $g=9.89\Umsq$，求抛掉的压舱物的质量。

%% answer:
\jdanswer{
$m'=101\Ukg$
}

%% memo:
\memoanswer{
由牛顿第二定律得：$mg-f=ma$，又 $h=v_{0}t+\frac{1}{2}at^{2}$，解得 $a=1\Umsq$。抛出压舱物后气球减速下降，有 $f-(m-m')g=(m-m')a'$，且 $\Delta v=a'\Delta t$，其中 $\frac{\Delta v}{\Delta t}=\frac{3}{300}\Umsq=0.01\Umsq$。解得：$m'=m\frac{a+\Delta v/\Delta t}{g+\Delta v/\Delta t}=990\times\frac{1+0.01}{9.89+0.01}\Ukg=101\Ukg$。
}

\item
%% number: 16
%% paperName: 2008年海南高考第 16 题 11 分
%% typeId: 6
%% type: 计算题
%% chapter: 磁场
%% point: 带电粒子在复合场中的运动
%% method: 无
%% score: 11
%% degree: 350
%% duplicateId: 0
%% body:
如图，空间存在匀强电场和匀强磁场，电场方向为 $y$ 轴正方向，磁场方向垂直于 $xy$ 平面（纸面）向外，电场和磁场都可以随意加上或撤除，重新加上的电场或磁场与撤除前的一样。一带正电荷的粒子从 $P(x=0,y=h)$ 点以一定的速度平行于 $x$ 轴正向入射。这时若只有磁场，粒子将做半径为 $R_{0}$ 的圆周运动；若同时存在电场和磁场，粒子恰好做直线运动。现在，只加电场，当粒子从 P 点运动到 $x=R_{0}$ 平面（图中虚线所示）时，立即撤除电场同时加上磁场，粒子继续运动，其轨迹与 $x$ 轴交于 M 点。不计重力。求：

\begin{enumerate}
	\item
	粒子到达 $x=R_{0}$ 平面时速度方向与 $x$ 轴的夹角以及粒子到 $x$ 轴的距离；
	\item
	M 点的横坐标 $x_{M}$。
\end{enumerate}

\onepicture[w=5.2cm, align=right]{figs/16.png}

%% answer:
\jdanswer{
\begin{enumerate}
	\item
	速度方向与 $x$ 轴夹角为 $\theta=\frac{\pi}{4}$；粒子与 $x$ 轴的距离为 $H=h+\frac{R_{0}}{2}$。
	
	\item
	$x_{M}=2R_{0}+\sqrt{\frac{7}{4}R_{0}^{2}+R_{0}h-h^{2}}$。
\end{enumerate}
}

%% memo:
\memoanswer{
（1）做直线运动有 $qE=qBv_{0}$，做圆周运动有 $qBv_{0}=m\frac{v_{0}^{2}}{R_{0}}$。只有电场时，粒子做类平抛运动，有 $qE=ma$，$R_{0}=v_{0}t$，$v_{y}=at$，解得 $v_{y}=v_{0}$。粒子速度大小为 $v=\sqrt{v_{0}^{2}+v_{y}^{2}}=\sqrt{2}v_{0}$，速度方向与 $x$ 轴夹角为 $\theta=\frac{\pi}{4}$，粒子与 $x$ 轴的距离为 $H=h+\frac{1}{2}at^{2}=h+\frac{R_{0}}{2}$。

（2）撤去电场加上磁场后，有 $qBv=m\frac{v^{2}}{R}$，解得 $R=\sqrt{2}R_{0}$。粒子运动轨迹如图所示，圆心 C 位于与速度 $v$ 方向垂直的直线上，该直线与 $x$ 轴和 $y$ 轴的夹角均为 $\pi/4$，由几何关系得 C 点坐标为 $x_{C}=2R_{0}$，$y_{C}=H-R_{0}=h-\frac{R_{0}}{2}$。过 C 作 $x$ 轴的垂线，在 $\Delta CDM$ 中：$\overline{CM}=R=\sqrt{2}R_{0}$，$\overline{CD}=y_{C}=h-\frac{R_{0}}{2}$，解得 $\overline{DM}=\sqrt{\overline{CM}^{2}-\overline{CD}^{2}}=\sqrt{\frac{7}{4}R_{0}^{2}+R_{0}h-h^{2}}$。M 点横坐标为 $x_{M}=2R_{0}+\sqrt{\frac{7}{4}R_{0}^{2}+R_{0}h-h^{2}}$。

\onepicture[w=5.0cm]{figs/16_an.jpg}
}

\item
%% number: 17
%% paperName: 2008年海南高考第 17 题 8 分
%% typeId: 6
%% type: 计算题
%% chapter: 气体性质
%% point: 气体多过程
%% method: 无
%% score: 8
%% degree: 550
%% duplicateId: 0
%% body:
\begin{enumerate}
	\item
	（4 分）下列关于分子运动和热现象的说法正确的是\tkanswer[1.5cm]{BCE}（填入正确选项前的字母，每选错一个扣 1 分，最低得分为 0 分）。
	
	\sixchoices[answer=BCE]
	{气体如果失去了容器的约束就会散开，这是因为气体分子之间存在势能的缘故}
	{一定量 $100^{\circ}$C 的水变成 $100^{\circ}$C 的水蒸汽，其分子之间的势能增加}
	{对于一定量的气体，如果压强不变，体积增大，那么它一定从外界吸热}
	{如果气体分子总数不变，而气体温度升高，气体分子的平均动能增大，因此压强必然增大}
	{一定量气体的内能等于其所有分子热运动动能和分子之间势能的总和}
	{如果气体温度升高，那么所有分子的速率都增加}
	
	\item
	（8 分）如图，一根粗细均匀、内壁光滑、竖直放置的玻璃管下端密封，上端封闭但留有一抽气孔。管内下部被活塞封住一定量的气体（可视为理想气体），气体温度为 $T_{1}$。开始时，将活塞上方的气体缓慢抽出，当活塞上方的压强达到 $p_{0}$ 时，活塞下方气体的体积为 $V_{1}$，活塞上方玻璃管的容积为 $2.6V_{1}$。活塞因重力而产生的压强为 $0.5p_{0}$。继续将活塞上方抽成真空并密封。整个抽气过程中管内气体温度始终保持不变。然后将密封的气体缓慢加热。求：①活塞刚碰到玻璃管顶部时气体的温度；②当气体温度达到 $1.8T_{1}$ 时气体的压强。
\end{enumerate}

\onepicture[w=2.2cm, align=right]{figs/17.png}

%% answer:
\jdanswer{
\begin{enumerate}
	\item
	$BCE$
	
	\item
	① $T'=1.2T_{1}$；② $p_{2}=0.75p_{0}$。
\end{enumerate}
}

%% memo:
\memoanswer{
（1）A 错误之处在于气体分子是无规则运动的，故失去容器后就会散开；D 选项中没考虑气体的体积对压强的影响；F 选项对温度升高，分子平均动能增大、平均速率增大，但不是每个分子的速率都增大，对单个分子的研究是毫无意义的。故选 BCE。

（2）①从活塞上方的压强达到 $p_{0}$ 到活塞上方抽成真空的过程为等温过程：$1.5p_{0}V_{1}=0.5p_{0}V_{2}$，得 $V_{2}=3V_{1}$；缓慢加热，当活塞刚碰到玻璃管顶部时为等压过程：$\frac{3V_{1}}{T_{1}}=\frac{3.6V_{1}}{T_{2}}$，得 $T_{2}=1.2T_{1}$，即 $T'=1.2T_{1}$。②继续加热到 $1.8T_{1}$ 时为等容过程：$\frac{0.5p_{0}}{1.2T_{1}}=\frac{p}{1.8T_{1}}$，解得 $p=0.75p_{0}$。
}

\item
%% number: 18
%% paperName: 2008年海南高考第 18 题 8 分
%% typeId: 6
%% type: 计算题
%% chapter: 机械波
%% point: 波长频率波速
%% method: 无
%% score: 8
%% degree: 550
%% duplicateId: 0
%% body:
\begin{enumerate}
	\item
	（4 分）设宇宙射线粒子的能量是其静止能量的 $k$ 倍。则粒子运动时的质量等于其静止质量的\tkanswer[1cm]{$k$}倍，粒子运动速度是光速的\tkanswer[2.5cm]{$\frac{\sqrt{k^{2}-1}}{k}$}倍。
	
	\item
	（8 分）某实验室中悬挂着一弹簧振子和一单摆，弹簧振子的弹簧和小球（球中间有孔）都套在固定的光滑竖直杆上。某次有感地震中观察到静止的振子开始振动 $4.0\Us$ 后，单摆才开始摆动。此次地震中同一震源产生的地震纵波和横波的波长分别为 $10\Ukm$ 和 $5.0\Ukm$，频率为 $1.0\UHz$。假设该实验室恰好位于震源的正上方，求震源离实验室的距离。
\end{enumerate}

%% answer:
\jdanswer{
\begin{enumerate}
	\item
	$k$；$\frac{\sqrt{k^{2}-1}}{k}$。
	
	\item
	$s=40\Ukm$。
\end{enumerate}
}

%% memo:
\memoanswer{
（1）以速度 $v$ 运动时的能量 $E=mc^{2}$，静止时的能量为 $E_{0}=m_{0}c^{2}$，依题意 $E=kE_{0}$，故 $m=km_{0}$；由 $m=\frac{m_{0}}{\sqrt{1-\frac{v^{2}}{c^{2}}}}$，解得 $v=\frac{\sqrt{k^{2}-1}}{k}c$，即粒子运动时的质量等于静止质量的 $k$ 倍，速度是光速的 $\frac{\sqrt{k^{2}-1}}{k}$ 倍。

（2）地震纵波传播速度为 $v_{P}=f\lambda_{P}$，地震横波传播速度为 $v_{S}=f\lambda_{S}$。设震源离实验室距离为 $s$，有 $s=v_{P}t$，$s=v_{S}(t+\Delta t)$，解得 $s=\frac{f\Delta t}{\frac{1}{\lambda_{S}}-\frac{1}{\lambda_{P}}}=\frac{1.0\times4.0}{\frac{1}{5.0}-\frac{1}{10}}\Ukm=40\Ukm$。
}

\item
%% number: 19
%% paperName: 2008年海南高考第 19 题 8 分
%% typeId: 6
%% type: 计算题
%% chapter: 运动和力
%% point: 动量和能量
%% method: 无
%% score: 8
%% degree: 550
%% duplicateId: 0
%% body:
\begin{enumerate}
	\item
	（4 分）某考古队发现一古生物骸骨。考古专家根据骸骨中 $\ce{^{14}_{6}C}$ 的含量推断出了该生物死亡的年代。已知此骸骨中 $\ce{^{14}_{6}C}$ 的含量为活着的生物体中 $\ce{^{14}_{6}C}$ 的 $1/4$，$\ce{^{14}_{6}C}$ 的半衰期为 5730 年。该生物死亡时距今约\tkanswer[2cm]{$1.1\times10^{4}$}年。
	
	\item
	（8 分）一置于桌面上质量为 $M$ 的玩具炮，水平发射质量为 $m$ 的炮弹。炮可在水平方向自由移动。当炮身上未放置其它重物时，炮弹可击中水平地面上的目标 A；当炮身上固定一质量为 $M_{0}$ 的重物时，在原发射位置沿同一方向发射的炮弹可击中水平地面上的目标 B。炮口离水平地面的高度为 $h$。如果两次发射时“火药”提供的机械能相等，求 B、A 两目标与炮弹发射点之间的水平距离之比。
\end{enumerate}

%% answer:
\jdanswer{
\begin{enumerate}
	\item
	$1.1\times10^{4}$（11460 或 $1.0\times10^{4}\sim1.2\times10^{4}$ 均可）。
	
	\item
	$\sqrt{\frac{(M+M_{0})(M+m)}{M(M+M_{0}+m)}}$。
\end{enumerate}
}

%% memo:
\memoanswer{
（1）该核剩下 $1/4$，说明正好经过两个半衰期时间，故该生物死亡时距今约 $2\times5730$ 年 $=11460$ 年。

（2）由动量守恒定律和能量守恒定律得：$0=mv_{1}-Mv_{2}$，$E=\frac{1}{2}mv_{1}^{2}+\frac{1}{2}Mv_{2}^{2}$，解得 $v_{1}=\sqrt{\frac{2EM}{m(M+m)}}$。炮弹射出后做平抛，有 $h=\frac{1}{2}gt^{2}$，$x=v_{1}t$，解得目标 A 距炮口的水平距离为 $x=\sqrt{\frac{4EMh}{gm(M+m)}}$。同理，目标 B 距炮口的水平距离为 $x'=\sqrt{\frac{4E(M+M_{0})h}{gm(M+M_{0}+m)}}$。解得 $\frac{x'}{x}=\sqrt{\frac{(M+M_{0})(M+m)}{M(M+M_{0}+m)}}$。
}

\end{enumerate}

\end{document}
